% @order: 1
% @description: 设置全局字体、页面布局(数模论文专用)
% @\usepackage{abstract}

% !字体和版式设置
% #region

% 页面布局
\geometry{
    left=2.5cm,
    right=2.5cm,
    top=3cm,
    bottom=3cm
}

% 中文字体设置
%\setCJKmainfont{SimSun}[BoldFont=SimHei,ItalicFont=KaiTi]
\setCJKmainfont{Source Han Serif SC}[
    BoldFont={Source Han Sans SC Medium},
    ItalicFont={Kaiti SC}
]
\setCJKsansfont{Source Han Sans SC}[
    BoldFont={Source Han Sans SC Bold}
]
\newCJKfontfamily\heavyfont{Source Han Serif SC Heavy}
\setCJKmonofont{PingFang SC}

% 数学字体设置
\setmathfont{latinmodern-math.otf}
\setmathfont[
    range={"2205},
    Scale=MatchUppercase
]{STIX Two Math}
\sisetup{
    per-mode = reciprocal,
    inter-unit-product = \ensuremath{\cdot}
}

% 文档分页与数学排版
\let\cleardoublepage\clearpage
\allowdisplaybreaks

% 页面垂直对齐
\raggedbottom

\ctexset{ % 设置一级标题格式
	section={
		format={\centering\Large\textbf},
		number = \chinese{section},
		}
	}

% 表格列样式设置
\renewcommand{\tabularxcolumn}[1]{m{#1}}
\newcolumntype{L}{>{\hbadness=10000\relax}X}
\newcolumntype{Y}{>{\RaggedRight\arraybackslash}X} % 左对齐可换行
\newcolumntype{Z}{>{\centering\arraybackslash}X}   % 居中对齐

% abstract环境设置
\ctexset{abstractname={}} % 取消默认摘要标题
\newcommand{\keywordscontent}{} 
\newcommand{\keywords}[1]{
    \def\keywordscontent{#1}
}
\renewenvironment{abstract}{
    \par
    \noindent\begin{center}{\Large\bfseries 摘\hspace{2em}要}\end{center}\par % 摘要标题
    \def\keywordscontent{} % 清空关键词
}{
    \par
    \ifx\keywordscontent\empty
    \else % 输出关键词
        \vspace{0.8em}
        {\noindent\bfseries 关键词：} \textit{\keywordscontent}\par
    \fi
}

% 符号说明表格样式设置
\newenvironment{symboltable}[1]{
    \begin{xltabular}{\textwidth}{
        @{} >{\hsize=0.2\hsize}Z >{\hsize=0.6\hsize}Z >{\hsize=0.2\hsize}Z @{}
    }
    \caption{符号说明及其单位} \label{tab:dataName} \\
    \toprule[0.08em]
    \textbf{符号} & \textbf{说明} & \textbf{单位} \\
    \midrule[0.05em]
    \endfirsthead % 第一页表头
    
    \multicolumn{3}{c}{表\thetable \quad 符号含义及其单位（续）} \\ 
    \toprule[0.08em]
    \textbf{符号} & \textbf{说明} & \textbf{单位} \\
    \midrule[0.05em]
    \endhead % 后续页表头
    
    \midrule[0.05em]
    \multicolumn{3}{r@{}}{（接下表）} \\
    \endfoot % 非尾页的页脚
    
    \bottomrule[0.08em]
    \endlastfoot % 最后一页底部
    
    #1 % 插入表格数据内容
    
    \end{xltabular}
}
% #endregion
